\section{Related Work}
\label{sec:related-work}

Research on concept drift in process mining has considered several ways to describe and compare process behavior over time. Control-flow approaches use relations between activities or distributions of execution traces. Other approaches track timing, resource use, and case characteristics. Changes can also be detected by monitoring conformance to process models or constraints, by recognizing patterns learned from labeled examples, or by grouping similar observations and tracking shifts in those groups.

\subsection{Activity Relations and Trace Distributions}

Bose et al. characterize control flow through relations between activities. Their relation type count records how many activities always, sometimes, or never follow each activity in a log fragment \cite{bose2011handling}. ProDrift tests whether the frequencies of partially ordered runs change across neighboring windows. By grouping traces that differ only in the order of concurrent activities, it avoids treating every observed ordering as a different run \cite{maaradji2017detecting}. The event-stream method of Ostovar et al. compares counts of relations extracted from events, allowing it to produce evidence while cases are still open \cite{ostovar2016streams}.

Local Completeness-based Drift Detection looks for relations that appear or cease to occur, while relation invariance uses changes in the presence of directly-follows and eventually-follows relations \cite{lin2022lcdd,zheng2017detecting}. Graph metrics summarize the structure of a process graph, for example through branching and connectivity \cite{seeliger2017graph}. In Kairo, activity and relation features help describe the observed prefix of a running case. Similar prefixes can share a state, and drift is assessed through changes in the frequencies of those states. The resource and inter-case analyses use other measurements and apply the same comparison to their state frequencies.

\subsection{Drift beyond Control Flow}

Brockhoff et al. use the Earth Mover's Distance to compare distributions of traces. The distance accounts for how similar traces are, and the time-aware representation includes service or sojourn times \cite{brockhoff2020time}. This retains differences between traces that can be lost when similar observations are assigned to the same state. Kairo can also represent conditions such as workload and resource allocation across cases in a calendar window.

Adams et al. detect and relate changes in feature series covering multiple process perspectives \cite{adams2023explainable}. Klijn et al. extend multi-perspective drift detection to actor behavior and collaboration using an event knowledge graph \cite{klijn2024multi}. Khraiwesh et al. use the coefficient of variation to choose a smoothing width for each indicator before detecting change points \cite{khraiwesh2027window}. Multivariate tests such as maximum mean discrepancy can detect changes in the joint distribution of feature vectors \cite{gretton2012kernel}. Kairo groups related measurements into state descriptions for each perspective, such as a workload state that combines arrivals with the cases already in progress. The resulting state can recur, and its share can be examined alongside the measurements typical of that state.

\subsection{Conformance, Constraints, and Supervised Detection}

FAIRY monitors how incoming traces conform to a reference model and updates the model after a detected change \cite{pasquadibisceglie2026fairy}. CRIER tracks fitness and precision over windows. Precision can help reveal a change when cases still fit the model but some allowed behavior disappears from the observations \cite{gallego2025crier}. Visual Drift Detection examines changes in declarative constraints, such as how regularly one activity is followed by another \cite{yeshchenko2019visual}. These methods provide evidence in terms of the model or rules they monitor. Kairo learns state descriptions directly from the selected log features and uses a reference distribution of states.

Kraus and van der Aa use labeled examples to train a detector to recognize drift patterns in images of window similarities \cite{kraus2025machine}. Kairo learns its states without drift annotations. Its divergence signal indicates when state usage changes and can guide inspection of the underlying observations.

\subsection{Clustering and Process States}

Clustering provides a way to summarize behavior before looking for change. Barbon Junior et al. cluster traces in an event stream and use the emergence of a new cluster as a signal \cite{barbonjunior2018framework}. Version Clustering groups windows into process versions based on similarities in their directly-follows profiles and estimates boundaries between versions \cite{abarcazuniga2026version}. Kairo uses clustering to summarize event prefixes or calendar windows, depending on the perspective. Several states can be present throughout a stable period, and a change may be visible as a shift in their relative frequency.

The state representation builds on the work of Berti et al., who use a self-organizing map to learn execution states from sequences of events associated with a business object \cite{berti2025ocel}. Kairo uses case prefixes to describe within-case behavior and extends the state representation to resource and inter-case conditions over calendar time. Steady-state detection by Kraus et al. identifies periods of stable process behavior for tasks such as discovery and performance analysis \cite{kraus2026steady}. Kairo uses states as recurring descriptions of observed behavior and tracks how their distribution changes over time.
